\section{The Rabung criterion and certificates}\label{s:method}

\subsection{Power-residue colourings}\label{ss:prc}

The certification argument has three layers.  Rabung's theorem turns a finite
criterion into a lower bound; Proposition~\ref{p:bstar} proves the equivalent
operational form used here; and a stand-alone, separate-source checker
re-establishes that criterion for each reported triple.  The GPU search of
Section~\ref{s:campaign} discovers candidate triples but is not a premise of
this implication.

An \emph{explicit certificate} for $W(r,k)>n$ is an $r$-colouring of
$\{1,\dots,n\}$ with no monochromatic $k$-term progression.  The large claims
instead use a \emph{succinct Rabung certificate} $(p,r,k)$, from which the
checker reconstructs the hypotheses of the construction theorem.  It does not
enumerate every progression in the resulting colouring.

Let $p$ be a prime with $p\equiv 1 \pmod r$ and $p>k$, let $g$ be a primitive
root modulo $p$, and let $\operatorname{ind}(x)$ denote the discrete logarithm
of $x$ to base $g$. The \emph{power-residue colouring} is
\[
  c(x) \;=\; \operatorname{ind}(x) \bmod r , \qquad x \in \{1,\dots,p-1\},
\]
i.e.\ the partition of $\mathbb{Z}_p^{*}$ into the $r$ classes of $r$-th power
residues. Its key property is multiplicativity: $c(xy)=c(x)+c(y) \pmod r$.
Consequently, dilating a progression $\{a,a+d,\dots,a+(k-1)d\}$ with
$p \nmid d$ by $u=d^{-1} \bmod p$ shifts all colours by the constant $c(u)$ and
maps the progression to $k$ \emph{consecutive} residues modulo $p$. The search
for monochromatic progressions therefore reduces to a linear scan of the run
lengths of the sequence $c(1),c(2),\dots,c(p-1)$, plus a boundary condition at
the multiples of $p$; this reduction from an $O(p^2/k)$ test to an $O(p)$ test
per prime is what makes a mass scan possible.

\subsection{Rabung's theorem and the operational criterion}\label{ss:criterion}

Rabung~\cite{rabung1979} proved that the extension $\vartheta'$ of the
power-residue colouring to $[0,(k-1)p]$, with the multiples of $p$ colored by any
non-constant assignment, is free of monochromatic $k$-term progressions
\emph{if and only if}
\begin{itemize}
\item[(a)] no $k$ consecutive positions in $[1,p-1]$ receive the same colour
  (max run $<k$); and
\item[(b)] a boundary condition holds. Put
  $\varepsilon:=c(-1)=c(p-1)\in\{0,r/2\}$, where the second value is possible
  only for even $r$. In Rabung's multiplicative notation the corresponding
  sign is $+1$ when $\varepsilon=0$ and $-1$ when $\varepsilon=r/2$. If
  $\varepsilon=0$ the integers
  $1,2,\dots,\lceil (k-1)/2 \rceil$ must not all lie in the same class, while if
  $\varepsilon=r/2$ the same is required of $1,2,\dots,k-1$. A \emph{singleton}
  counts as ``all in the same class'', and so fails~(b).
\end{itemize}
If $r$ is odd, $(p-1)/r$ is even and hence $\varepsilon=0$ necessarily.

\begin{remark}\label{r:ceiling}
Both details of~(b) matter. Rabung's bracket here denotes the ceiling.
Replacing it by a floor, or treating a
singleton as vacuously satisfying the condition, yields a criterion that is
\emph{not} equivalent to Rabung's.  For $p=5$, $r=2$, $k=4$, condition~(a) holds and the ceiling
correctly accepts the boundary test, whereas the floor rejects it. For
$p=5$, $r=2$, $k=3$, treating the singleton as acceptable gives a false
acceptance of the complete criterion.  The ceiling and floor coincide only for odd~$k$.
Both failure modes occurred during development and were detected by comparison
with the finite oracle described below.
\end{remark}
Since $[0,(k-1)p]$ contains $(k-1)p+1$ integers, a prime satisfying
(a)$\wedge$(b) yields
\begin{equation}\label{eq:bound}
  W(r,k) \;>\; (k-1)\,p+1 .
\end{equation}
We call such a prime \emph{valid} for $(r,k)$.

Our implementations test an equivalent form of (b), denoted (B$^{*}$), which is
uniform in the parity of $k$ and the value of $\varepsilon$: for every shift
$j\in\{0,\dots,k-1\}$, the $k-1$ character positions
$\{-j,\dots,-1,1,\dots,k-1-j\}$ of a $k$-window centered on a multiple of $p$
must not all receive the same colour. The equivalence rests on the
multiplicative structure: if a single shifted window is monochromatic, its
dilates realize every colour at the central position (the congruence
$p\equiv1\pmod r$ makes $c$ a surjective homomorphism to $\mathbb Z/r\mathbb Z$), so the joker position cannot be
colored at all; otherwise any non-constant colouring of the multiples of $p$
works.  Proposition~\ref{p:bstar} gives the mathematical proof; the finite
comparison below audits the implementation rather than replacing that proof.

\paragraph{A small example.}
Take $(p,r,k)=(5,2,4)$ and primitive root $g=2$. The nonzero residue colours
are $(c(1),c(2),c(3),c(4))=(0,1,1,0)$. The longest run is $2<4$, and the
boundary interval $\{1,2\}$ contains both colours. Thus the criterion gives
$W(2,4)>16$: the colouring on $[0,15]$ consists of three repeated nonzero
residue blocks separated by four joker positions. Every non-constant
assignment to those jokers works. This example illustrates the construction;
it is not asserted to be the best bound at this small length. It also shows
why replacing the ceiling by a floor would discard a valid certificate.

\begin{lemma}[Monotonicity at fixed prime]\label{lem:k-monotone}
Let $r\ge2$, let $3\le k\le \ell<p$, and let $p\equiv1\pmod r$ be prime. If
$p$ satisfies $(a)\wedge(B^{*})$ for $(r,k)$, then it satisfies the same
criterion for $(r,\ell)$. Consequently
\[
  W(r,\ell)>(\ell-1)p+1.
\]
\end{lemma}
\begin{proof}
It is enough to pass from $k$ to $k+1$. A monochromatic run of length $k+1$
would contain one of length $k$, so condition~(a) persists. Write
\[
 S_j^{(m)}=\{-j,\dots,-1\}\cup\{1,\dots,m-1-j\},
 \qquad 0\le j\le m-1.
\]
For $0\le j\le k-1$ we have $S_j^{(k)}\subset S_j^{(k+1)}$, and
$S_{k-1}^{(k)}\subset S_k^{(k+1)}$. Thus every boundary string for length
$k+1$ contains one already known to be non-monochromatic, so $(B^{*})$
persists. Iteration gives the claim.
\end{proof}

\appendix
\section{Equivalence of the operational criterion}\label{s:bstar}

The scan tests the criterion $(a)\wedge(B^{*})$ of Section~\ref{s:method} rather
than Rabung's $(a)\wedge(b)$ verbatim. The two are equivalent, and since every
bound in this paper rests on that equivalence we prove it here rather than rely
on the exhaustive machine comparison of Section~\ref{s:method}, which we regard
as a check on the \emph{implementation}, not as a substitute for a proof.

\paragraph{Setting.}
Let $p\equiv1\pmod r$ with $p>k\ge3$, let $g$ be a primitive root, and colour $x\in[1,p-1]$ by
$c(x)=\operatorname{ind}_g(x)\bmod r$; this is the power-residue colouring, and
$c(xy)=c(x)+c(y)$. Extend it to $[0,(k-1)p]$ by reducing modulo $p$ and colouring
the multiples of $p$ (the \emph{jokers}) by an arbitrary map
$\vartheta':\{0,p,\dots,(k-1)p\}\to\mathbb{Z}_r$. Write $\vartheta(-1)=c(p-1)$,
which is $0$ if $-1$ is an $r$-th power residue and $r/2$ otherwise (the latter
requires $r$ even).

An arithmetic progression of length $k$ and common difference $d$ is monochromatic
iff its terms all receive one colour. Progressions avoiding the jokers are handled
by~(a): by complete multiplicativity of $c$, the progression $a,a+d,\dots,a+(k-1)d$
reduces, after dividing by $d$ (which permutes the classes by the constant
$c(d)$), to the window $\{a/d,\,a/d+1,\dots,a/d+k-1\}$ of $k$ consecutive
residues. Hence \emph{some} joker-free progression is monochromatic iff \emph{some}
window of $k$ consecutive residues in $[1,p-1]$ is, which is exactly the negation
of~(a).

\paragraph{The joker condition.}
It remains to treat the progressions that meet a joker. First bound the common
difference: a $k$-term progression inside $[0,(k-1)p]$ has
$(k-1)d\le(k-1)p$, so $1\le d\le p$, and the case $d=p$ occurs for exactly one
progression, namely $0,p,2p,\dots,(k-1)p$, which consists \emph{entirely} of
jokers. That progression is monochromatic if and only if $\vartheta'$ is
constant---it is precisely the reason Rabung's statement requires a
\emph{non-constant} joker colouring, and any non-constant $\vartheta'$ breaks
it. For the remaining case $0<d<p$: since the progression has $k<p$ terms, two
of them congruent to $0$ modulo $p$ would force $p\mid(i'-i)d$ with
$0<|i'-i|<k<p$ and $p\nmid d$, which is impossible; so the progression meets at
most one multiple of $p$. Say it is the $(j+1)$-st term, $0\le j\le k-1$.
Dividing by $d$ as before, the progression becomes the \emph{shifted string}
\[
  S_j \;=\; \{-j,\dots,-1\}\cup\{1,\dots,k-1-j\}\pmod p,
\]
centred on the joker, together with the joker itself.

\begin{proposition}\label{p:bstar}
Fix $p\equiv1\pmod r$ and $k$. The colouring above is free of monochromatic
$k$-term progressions for \emph{some} choice of $\vartheta'$ if and only if
$(a)\wedge(B^{*})$ holds, where
\[
  (B^{*}):\qquad \text{no shifted string }S_j\ (0\le j\le k-1)\text{ is monochromatic.}
\]
Moreover, when $(a)\wedge(B^{*})$ holds, \emph{every} non-constant $\vartheta'$
works; and when it fails, \emph{no} $\vartheta'$ works. In particular the choice of
$\vartheta'$ is immaterial, as asserted by Rabung.
\end{proposition}

\begin{proof}
Suppose some $S_j$ is monochromatic, of colour $w$ say.

\emph{The lift.} We first check that every residue class of $d$ is realised by a
genuine progression inside $[0,(k-1)p]$ through a joker. Fix a representative
$d\in[1,p-1]$ and take the joker to be $jp$ and the first term to be
\[
  a \;=\; j\,(p-d).
\]
Then $a+jd=jp$, so the $(j+1)$-st term is the joker; the first term
$a=j(p-d)\ge0$ since $d<p$; and the last term is
\[
  a+(k-1)d \;=\; jp+(k-1-j)d \;\le\; jp+(k-1-j)(p-1) \;\le\; (k-1)p ,
\]
with equality in the last step only at $j=k-1$, where that last term is the
joker $jp$ itself; so all $k$ terms lie in $[0,(k-1)p]$. The non-joker terms are $jp+id$ for
$i\in S_j$, whose residues are $id \bmod p$; hence their colours are
$c(i)+c(d)=w+c(d)$.

\emph{The all-or-nothing mechanism.} As $d$ ranges over $(\mathbb{Z}/p)^{*}$, the
value $w+c(d)$ ranges over \emph{all} of $\mathbb{Z}_r$ (because $c$ is onto), and
by the lift each such value is attained by an admissible progression through the
\emph{same, fixed} joker $jp$. Hence for each colour $u\in\mathbb{Z}_r$ there is a
progression through $jp$ whose other $k-1$ terms are all of colour $u$; assigning
$\vartheta'(jp)=u$ makes it monochromatic. Every value of $\vartheta'$ at the
single joker $jp$ is thus forbidden, so no choice of $\vartheta'$---constant or
not---avoids a monochromatic progression. Conversely, if no $S_j$ is
monochromatic then no progression with $d<p$ through a joker has all its other
$k-1$ terms of a single colour, so no assignment of $\vartheta'$ can complete
one; the only remaining progression is the all-joker one ($d=p$), which any
\emph{non-constant} $\vartheta'$ breaks; combined with~(a) for the joker-free
progressions, every non-constant $\vartheta'$ works.  The $d=p$ case is exactly
what consumes the non-constancy hypothesis.  Thus
$(a)\wedge(B^{*})$ is all-or-nothing over non-constant $\vartheta'$, which is
why Rabung may state it without reference to the joker colouring.
\end{proof}

\paragraph{Reduction of $(B^{*})$ to a single interval.}
Since $c(-x)=c(x)+\vartheta(-1)$, the string $S_j$ is monochromatic iff the two
intervals $\{1,\dots,j\}$ and $\{1,\dots,k-1-j\}$ are each monochromatic and their
colours agree after the shift by $\vartheta(-1)$.

\emph{Case $\vartheta(-1)=0$.} Then $c(-x)=c(x)$, so $S_j$ is monochromatic iff the
interval $\{1,\dots,\max(j,\,k-1-j)\}$ is. Minimising $\max(j,k-1-j)$ over
$0\le j\le k-1$ gives $\lceil(k-1)/2\rceil$, attained at $j=\lfloor(k-1)/2\rfloor$.
Hence
\[
  (B^{*})\iff \{1,2,\dots,\lceil (k-1)/2\rceil\}\ \text{is not monochromatic.}
\]
A one-element interval \emph{is} monochromatic, and therefore fails the condition;
under the standing assumption $k\ge3$ this is the case $k=3$, and it is the reason
a singleton must be counted as ``all in the same class''.

\emph{Case $\vartheta(-1)=r/2$ ($r$ even).} Now $c(-x)=c(x)+r/2\ne c(x)$, so a
two-sided string can never be monochromatic: only the one-sided strings $j=0$ and
$j=k-1$ survive, both giving the interval $\{1,\dots,k-1\}$. Hence
\[
  (B^{*})\iff \{1,2,\dots,k-1\}\ \text{is not monochromatic,}
\]
Unlike the length-$k$ prohibition of~(a), this is a condition on an interval of
length $k-1$, and is therefore \emph{not} implied by~(a): it is a genuine
additional constraint, and on a longer interval than in the previous case.

\begin{remark}
Both refinements matter. Replacing $\lceil\cdot\rceil$ by $\lfloor\cdot\rfloor$, or
treating a singleton as vacuously satisfying the condition, gives a boundary
condition that is not equivalent to~(b).  Counterexamples occur already at
$p=13$, $r=2$, $k=4$ for the floor variant and at $p=5$, $r=2$, $k=3$ for the
singleton convention; condition~(a) can mask some such boundary errors.
\end{remark}

This proposition is the mathematical bridge used for every one of the thirteen
registered triples: once a verifier establishes $(a)\wedge(B^{*})$, Rabung's
theorem yields the corresponding direct inequality.


\subsection{Finite oracle audit of the implementation}\label{ss:oracle}

We tested the implementation of (a)$\wedge$(B$^{*}$) against a frozen reference
oracle that explicitly builds the colouring of $[1,(k-1)p+1]$ and exhaustively
checks all $k$-term progressions. One bookkeeping detail matters: Rabung's
construction lives on $[0,(k-1)p]$, and the oracle works on its translate by
one---it colours $x$ by the class of $(x-1)\bmod p$ and places the jokers at
$x\equiv1\pmod p$---so the two objects are literally the same colouring up to
relabeling of positions, and freedom from monochromatic progressions is
translation-invariant. A sweep over all primes $p<10^{4}$ with
$p\equiv 1 \pmod r$, $p>k$, $r\in[2,6]$, $k\in[3,8]$ --- $20\,170$ cases, of
which $4\,499$ were valid according to the criterion --- reported zero
disagreements.  A fresh staging replay reproduced those counts and its exact
source identity, environment, and output streams are archived with this
repository. This larger oracle sweep is historical evidence; the final CPU
suite separately repeats the small oracle cases and checks the retained record. A
stronger test enumerates all $r^{k}$ colourings of the joker positions
($p<150$, $r\le 4$, $k\le 5$) and confirms both directions of the equivalence:
(a)$\wedge$(b) holds if and only if some valid extension exists, and when it
holds every non-constant joker colouring is valid. We emphasize the epistemic
status: Rabung's theorem is proved in~\cite{rabung1979}; the agreement of our
operational form with the explicit oracle is exhaustively machine-checked on
the stated domain.

\subsection{Certificates and verification regimes}\label{ss:certificates}

For the bounds announced in Section~\ref{s:bounds} the certificate length
$(k-1)p+1$ is of order $3\times 10^{10}$ to $10^{11}$, so running the naive
quadratic oracle on the explicit certificate is not feasible. Our verification
policy therefore distinguishes two regimes. On the finite validation domain
actually reported above ($p<10^4$, $k\le8$, hence certificate length below
$7\times10^4$), the explicit certificate was checked directly by the reference
oracle.  We make no general feasibility claim for length $10^7$: a successful
naive oracle must inspect
\[
 \sum_{d=1}^{\lfloor(n-1)/(k-1)\rfloor}
       \bigl(n-(k-1)d\bigr)
\]
progressions. For the large primes, the proof of~\eqref{eq:bound} is Rabung's
theorem together with conditions (a)$\wedge$(b), and \emph{these} are what we
verify redundantly: the mirror, full-interval, and CPU-walk paths, followed by
separate CPU implementations of the full criterion.
Section~\ref{ss:validation} records which paths share source code or
arithmetic; cross-agreement is not described as statistical or implementation
independence.

At the large sizes of Section~\ref{s:bounds}, distributing the full colouring
would be neither feasible nor useful.  The reproducibility package therefore
distributes the triple $(p,r,k)$, from which the stand-alone, separate-source
verifier \texttt{verify\_claim} re-derives every required hypothesis.  Its
protocol is:
\begin{enumerate}
\item check that $p$ is prime by deterministic Miller--Rabin on the first
twelve prime bases $2,3,5,\dots,37$.  Sorenson and Webster's Theorem~1.1 gives
\[
  \psi_{12}=318\,665\,857\,834\,031\,151\,167\,461>2^{64},
\]
so this is exact through the unsigned $64$-bit
range~\cite{SorensonWebster2017}, far beyond the $p<2^{32}$ actually used;
\item check the preconditions $p\equiv1\pmod r$ and $p>k$;
\item \emph{find} a primitive root $g$ by trial, certifying primitivity against
the full trial-division factorization of $p-1$ (nothing about $g$ is taken as
input);
\item build the residue classes by modular exponentiation---$x^{(p-1)/r}$
compared against the powers of $\zeta=g^{(p-1)/r}$---using neither the
discrete-logarithm walk nor the GPU pipeline;
\item test (a)$\wedge$(B$^{*}$) and print \textsc{accept} or \textsc{reject}.
\end{enumerate}
For the fixed small colour counts used here, the checker takes
$O(p\log p)$ modular-arithmetic steps and $O(p)$ memory. It stores one byte
per residue at $r=2,3$: the largest certificate therefore requires about
$3.48$ GB for the colour array alone, before auxiliary memory. This is the
main practical cost of re-evaluating a large triple. The archived executions
already establish the finite checks used here; validating their identities
and rebuilding the paper are much lighter operations. No standardized
runtime benchmark is claimed. A streaming implementation could reduce memory,
but it is not the implementation whose executions are reported.
